\begin{theorem}[A direct upward crossing in the four-endpoint presentation]
    \label{thm:peps_first_string_physical_crossing}
    \lean{TNLean.PEPS.torusFirstStringCrossingBase,TNLean.PEPS.torusFirstStringCrossingGauge,TNLean.PEPS.torusFirstStringCrossingOutput,TNLean.PEPS.torusFirstStringCrossing_initial_extension,TNLean.PEPS.torusFirstStringCrossingOutput_eq_of_ne,TNLean.PEPS.torusFirstStringCrossingOutput_middle_transport,TNLean.PEPS.IsGIsometric.exists_unitary_torusFirstStringCrossing}
    \leanok
    Consider the native four-endpoint string assignment on a torus with
    $w\ge7$ and $H\ge6$. The upward movement tile based at $v+(1,1)$ has a
    common reconstruction gauge equal to $1$ on its lower row and $g$ on its
    other two rows. Its lower canonical insertion is $h^{-1}$. Extending this
    insertion to the middle chord changes only the bond based at $v+(1,2)$,
    whose rightward transport becomes $gh^{-1}g^{-1}$. The lower bond retains
    the continuing partner string, and all other ordered operators are retained.
    For an original regular $G$-isometric site tensor there exists one unitary
    on these six original spins, chosen before $g,h$ and every boundary label,
    which maps all actual input columns to the output columns. Its extension
    by the identity maps the actual closed input state to the output state.
    This is a direct elementary crossing auxiliary to
    \cite[Theorem~6.16 and Section~6.6]{Schuch2010PEPS}; it is not identified
    with the entire prescribed braid.
\end{theorem}
\begin{proof}
    \leanok
    \uses{thm:peps_two_plaquette_transport_table,thm:peps_torus_gauged_vertical_flux_move,thm:peps_internal_gauge_transport}
    The seven literal native transports determine the common gauge
    reconstruction. The canonical single-to-double chord movement then applies
    with the same gauge and exterior operators on both sides. Its physical
    unitary acts on every boundary column, hence on the closed contraction.
\end{proof}
